% Passes 5792--5799: matrix-ring affine model and transpose outer involution.
\subsection{The q=5 Latin carrier as an affine $M_2(\mathbb F_2)$ left--right geometry}
\label{sec:pass5792-5799-matrix-ring-transpose}

The Klein-$V_4$ chart reconstructed in Passes~5776--5783 admits a concrete
matrix-algebra model that resolves the source outer involution.  Choose
$V=W=\mathbb F_2^2$.  The full centreless order-$576$ Latin carrier is
\[
 \boxed{
 G_{576}=\operatorname{Hom}(W,V)\rtimes
 \bigl(GL(V)\times GL(W)\bigr).
 }
\]
After choosing bases,
$\operatorname{Hom}(W,V)\cong M_2(\mathbb F_2)$, so
\[
 \boxed{
 G_{576}\cong
 M_2(\mathbb F_2)_{+}:\bigl(GL_2(2)\times GL_2(2)\bigr),
 \qquad16\cdot6\cdot6=576.
 }
\]
Explicitly,
\[
 (X,A,B):M\longmapsto AMB^{-1}+X.
\]
The verifier enumerates all $576$ such affine maps.  The normal matrix space
contains six invertible and ten singular matrices, the classical
$M_2(\mathbb F_2)$ census.

The two degree-$12$ carriers have a particularly simple dual description.  The
point carrier is
\[
 \boxed{P=\{(w,x):0\ne w\in W,\ x\in V\}.}
\]
For $0\ne\phi\in V^*$ and $\psi\in W^*$ define
\[
 \boxed{
 H_{\phi,\psi}=\{(w,x)\in P:\phi(x)+\psi(w)=1\}.
 }
\]
There are twelve such six-subsets.  Under the explicit q=5 coordinate chart of
Pass~5778 they agree \emph{exactly} with the twelve frozen multiplicity-six
heavy supports.  Replacing the right side by zero gives exactly the twelve
Klein intercalate supports.  Hence the preceding complement theorem is the
single affine relation
\[
 \boxed{
 \phi(x)+\psi(w)=0/1
 \quad\Longleftrightarrow\quad
 \text{intercalate/heavy}.
 }
\]

The induced point action is
\[
 (w,x)\longmapsto\bigl(Bw,Ax+X(Bw)\bigr),
\]
while the dual heavy action is
\[
 (\phi,\psi)\longmapsto
 \left(\phi A^{-1},\psi B^{-1}+\phi A^{-1}X\right).
\]
The sixteen Reye/Latin lines are the graphs of the sixteen linear maps
$W\to V$ restricted to $W\setminus\{0\}$.  Exhausting all group elements gives
faithful order-$576$ actions on the point, heavy and line carriers, with the
two $12$-actions having stabilizer order $48$ and subdegrees $1,3,8$.  The
permutation-character Gram is
\[
 \boxed{
 \left(\langle\pi_i,\pi_j\rangle\right)_{i,j=P,H,L}
 =\begin{pmatrix}3&2&2\\2&3&2\\2&2&3\end{pmatrix},
 }
\]
and the point-side sign twist satisfies
\[
 \boxed{
 \bigl(
 \langle\varepsilon\pi_P,\pi_P\rangle,
 \langle\varepsilon\pi_P,\pi_H\rangle,
 \langle\varepsilon\pi_P,\pi_L\rangle
 \bigr)=(0,0,0).
 }
\]
Thus the independent affine-matrix model exactly replays the Pass~5667--5674
carrier-character certificate.

Now let $\Theta(M)=M^T$.  Direct conjugation gives
\[
 \boxed{
 \Theta\,(X,A,B)\,\Theta
 =\tau(X,A,B)
 =(X^T,B^{-T},A^{-T}),
 }
\]
with $\tau^2=1$.  Transpose normalizes the affine group but is not one of its
$576$ elements.  It fixes eight of the sixteen matrices.  More importantly,
with the chosen bases the bijection
\[
 \boxed{F:P\to H,\qquad F(w,x)=(w^T,x^T)}
\]
satisfies the exact intertwining law
\[
 \boxed{F(g\cdot p)=\tau(g)\cdot F(p)}
\]
for every group element and every point.  Hence the heavy action is precisely
the point action twisted by the transpose factor-swap automorphism.

The outerness is visible inside the normal translation subgroup
$T=M_2(\mathbb F_2)_+$.  A point stabilizer meets $T$ in
\[
 T_P=\{X:Xw=0\},
\]
a four-element fixed-right-kernel subspace whose aggregate image spans all of
$V$.  A heavy stabilizer meets $T$ in
\[
 T_H=\{X:\phi X=0\},
\]
a four-element fixed-left-image subspace whose aggregate image has dimension
one.  Internal conjugation acts on $T$ by $X\mapsto AXB^{-1}$ and cannot
exchange these two types, whereas
\[
 \boxed{T_P^T=T_H.}
\]
This supplies a concrete finite-algebra mechanism for the abstract carrier
outer involution of Pass~5674.

\paragraph{Prior art and evidence boundary.}
The ring $M_2(\mathbb F_2)$ and its use in two-qubit projective-ring geometry
are external prior art; see Saniga--Planat--Pracna,
\emph{Projective Ring Line Encompassing Two-Qubits}, arXiv:quant-ph/0611063.
That work uses the projective line over the noncommutative order-$16$ ring to
model the two-qubit Pauli commutation geometry.  The theorem here is narrower:
the q=5 Reye/Latin normal $16$-space is the \emph{additive} matrix space, its
order-$576$ carrier is the affine left/right unit action, and transpose
realizes the point/heavy factor swap.  No identification of q=5 W33 carrier
points with two-qubit observables, projective-ring points, entanglement classes,
or physical states is made.  No dynamics, gauge interaction, mass/coupling, or
physical-unification conclusion follows.
